Este trabalho parte de uma constatação sobre a forma como a informação aeronáutica chega aos seus usuários. Os produtos oficiais, como a AIP e as cartas aeronáuticas, apresentam o estado dos dados na data de efetividade de cada edição. Quando um valor muda, porém, não registram a causa da mudança, quem a afirmou nem as premissas de que o novo valor depende. No Aeroporto de São José dos Campos, a Carta de Obstáculos de 2022 e a Carta de Aeródromo de 2026 apresentam valores diferentes para o comprimento da pista, a zona de parada, a faixa de pista e a ASDA da cabeceira~16. Essas diferenças não constituem, por si sós, erros, já que as cartas têm propósitos distintos e foram emitidas com quatro anos de intervalo. Nenhuma das duas, contudo, permite identificar o que as motivou. O problema tratado é, portanto, de rastreabilidade, e não de exatidão.

Para esse problema, o trabalho propõe e avalia uma arquitetura de dados baseada em \textit{event sourcing}, inspirada no princípio do livro-razão contábil. Cada alteração é registrada como um evento acrescentado a um registro que não admite edição nem exclusão, com autor e causa declarados, e os eventos são encadeados por \textit{hash}. As posições são guardadas em relação a referentes, segundo o referenciamento linear da ISO~19148. Os valores que podem ser calculados, como comprimentos, azimutes e distâncias declaradas, são derivados das premissas no momento da consulta e nunca armazenados. A leitura é separada da escrita. O registro é publicado em ordem para consumidores que verificam a cadeia e montam os seus próprios modelos de leitura, e é entregue, ao final, como um banco extraído em arquivo único.

A arquitetura foi implementada como biblioteca de código aberto e avaliada segundo a pesquisa em ciência do design, por falseamento, em seis experimentos com dados reais da pista do SBSJ. Os critérios de aprovação foram definidos antes da execução. Os cinco experimentos sobre a arquitetura foram aprovados, e o primeiro, de caracterização, confirmou a lacuna (Tabela~\ref{tab:res-sintese}). A caracterização do E1 mostra que o modelo que guarda apenas os valores publicados não registra a causa de nenhum dos doze campos que mudam entre as cartas. Na arquitetura proposta, o comprimento recalculado a partir das coordenadas das cabeceiras, de 2675,342~m, coincide com o publicado dentro da precisão de publicação, o que mostra que o valor é reproduzível a partir das premissas, e não que o levantamento seja exato, e uma remedição altera somente o levantamento e o referencial. O estado de qualquer instante é reconstruído com a causa de cada mudança, e as alterações no meio da cadeia são detectadas, tanto no banco quanto no canal de distribuição. Os produtos gerados do mesmo registro não divergem, o que mostra consistência, e não correção da regra, e um verificador escrito apenas a partir da especificação pública chega ao mesmo \textit{hash} sem recorrer à biblioteca.

A resposta à pergunta de pesquisa é, assim, afirmativa no caso estudado, com limites bem localizados. Na pista avaliada, a arquitetura torna os dados rastreáveis até a causa, recalculáveis a partir das premissas, reconstruíveis no tempo e consistentes entre produtos, e as alterações no meio do histórico deixam evidência verificável. A auditabilidade obtida é documental. Ela garante que se saiba o que foi afirmado, por quem, quando e com base em quê, mas não que o afirmado seja verdadeiro, e o histórico reconstruído é o das afirmações registradas, e não o da condição física da pista: um fato descoberto depois aparece na data em que foi registrado. Dois limites dependem de elementos externos ao registro. A reescrita completa da cadeia só é detectada com uma âncora guardada fora do sistema, que a implementação permite, ao expor a cabeça da cadeia, mas não publica nem compara. A autoria é registrada como declarada, mas não comprovada, por falta de assinatura digital. Os demais limites decorrem de extensões ainda não implementadas e do escopo da avaliação, restrita a uma pista de um único aeródromo.

A contribuição do trabalho é a composição de mecanismos conhecidos, oriundos dos bancos de dados temporais, dos sistemas financeiros e dos registros verificáveis, para os dados de infraestrutura aeroportuária. Essa composição se acompanha das adaptações que o domínio exige, como as posições relativas a referentes, os derivados por regras normativas e a atomicidade imposta na escrita. A contribuição materializa-se em quatro artefatos, a saber, a arquitetura, a biblioteca, o banco extraído do SBSJ, publicado como conjunto de dados aberto, e um método de avaliação reprodutível. A arquitetura não substitui o AIXM nem o ciclo AIRAC. Ela os complementa, podendo produzir dados em AIXM como mais um modelo de leitura e oferecendo à emenda AIRAC uma âncora natural para o histórico.

\section{Trabalhos futuros}\label{sec:trabalhos-futuros}

Os trabalhos futuros mais imediatos atacam os limites que dependem de elementos externos. O primeiro é a âncora externa. Publicar a cabeça da cadeia junto com cada emenda AIRAC, ou depositá-la com o banco extraído em repositório com identificador persistente, permitiria detectar a reescrita completa do registro e as alterações no fim do canal de distribuição. O segundo é a assinatura digital das afirmações, que converteria o autor declarado em autor comprovado. Ambos podem ser incorporados sem alterar a forma do registro.

Em seguida, vêm as extensões já desenhadas. A norma de preenchimento em três camadas estenderia a todos os campos a distinção entre ``não informado'', ``não se aplica'' e ``desconhecido'', que hoje vale apenas para as premissas das distâncias declaradas. A extensão bitemporal acrescentaria o tempo de validade, declarado com autor e causa, e permitiria consultar o que era verdadeiro no mundo, e não apenas o que estava registrado. A escrita a partir de várias fontes, com fila de escrita e controle otimista por \textit{feature}, substituiria o bloqueio pessimista da implementação atual. A regra de distâncias declaradas para pistas com barreira de contenção completaria o conjunto de derivados.

A avaliação também pode ser ampliada. Um segundo estudo de caso, com várias pistas e outros elementos do lado ar, testaria a composição que hoje é apenas argumentada. Experimentos de desempenho mediriam o custo de leitura, o efeito dos caches e o comportamento do tópico sobre um sistema de mensagens real. Um experimento específico poderia verificar que apagar todas as cópias de leitura e reconstruí-las a partir do registro produz resultado idêntico. A interoperabilidade entre organizações exigiria, por sua vez, mais de um publicador.

No plano da aplicação, a extensão a dados operacionais é a direção de maior interesse para a engenharia aeronáutica. Relatos de ocorrência e eventos de voo registrados como \textit{features} permitiriam relacionar a operação ao estado da infraestrutura no instante de cada ocorrência e estudar o efeito de mudanças na infraestrutura sobre a operação. A emissão de dados em AIXM e de NOTAM digitais a partir do registro, bem como uma interface de consulta sobre as projeções dos consumidores, aproximariam a arquitetura do uso corrente. Ferramentas de consulta para agentes de inteligência artificial poderiam, ainda, aproveitar a proveniência citável e a consulta no tempo já demonstradas.

O ponto de partida deste trabalho foi uma pergunta simples, a saber, por que um valor publicado mudou. As publicações aeronáuticas, tal como chegam hoje aos usuários, não permitem respondê-la. A arquitetura proposta mostra, no caso estudado, que é possível conservar essa resposta junto com o dado, de forma verificável por qualquer pessoa e com tecnologia comum. Trata-se de um passo para que a rastreabilidade exigida da cadeia de informação aeronáutica chegue também a quem utiliza os dados.